\section{System Architecture}\label{sec:architecture}

\subsection{Processes}
\Cref{fig:arch} shows two users and the tracker. Every user runs a \emph{peer daemon} and talks to it
through a browser; the browser never talks to the tracker or to another peer. The daemons talk to the
tracker over HTTP, to each other over TCP for chat and over UDP for files.

\tikzfig[fig:arch]{arch}{Architecture: two peer daemons, their browsers and the tracker. The
tracker can also relay an encrypted chat frame when a peer cannot reach another one directly.}

\begin{table}[H]\centering\footnotesize
\caption{Processes and their ports (defaults of the README)}\label{tab:processes}
\begin{tabularx}{\textwidth}{@{}l l Y@{}}
\toprule
\textbf{Process} & \textbf{Ports} & \textbf{Role} \\
\midrule
Tracker (\code{start_tracker.py}) & HTTP 8001 & accounts, registry of live peers, relay fallback, status page \\
Peer daemon (\code{start_sampleapp.py}) & HTTP 8000, TCP 9101, UDP 9102 & web UI and REST API of one user; P2P chat listener; UDP file endpoint \\
Reverse proxy (\code{start_proxy.py}) & HTTP 8080 & optional: virtual hosts, round-robin / random / least-connections, fail-over \\
Static backend (\code{start_backend.py}) & HTTP 9000 & optional: serves \file{www/} and \file{static/}; target for the proxy \\
\bottomrule
\end{tabularx}
\end{table}

\subsection{Why a hybrid design}
A pure client-server chat would route every message through one machine, which then sees all the
traffic and is a single point of failure. A pure peer-to-peer design has no natural place for the
user database or for discovery. The hybrid design keeps the small, rarely changing state --- who is
online and where --- on the tracker, and moves the traffic itself between the peers. As a result the
tracker can fail without stopping a running conversation: peers that are already connected keep
chatting with the last peer list they received (tested in \cref{sec:testing}).

\subsection{Separation of concerns}
The assignment requires network logic, data processing and the user interface to be separated
\cite{brief1}. \Cref{tab:modules} shows how the code is divided; no module reaches across two layers.

\begin{table}[H]\centering\footnotesize
\caption{Modules by layer}\label{tab:modules}
\begin{tabularx}{\textwidth}{@{}l >{\raggedright\arraybackslash}p{4.6cm} Y@{}}
\toprule
\textbf{Layer} & \textbf{Module} & \textbf{Responsibility} \\
\midrule
User interface & \file{static/chat.html}, \file{app.js}, \file{chat.css} & rendering, polling, user actions; client-side only \\
HTTP API & \file{apps/chat_app.py}, \file{apps/tracker_app.py} & routes, input validation, sessions, mapping errors to status codes \\
Application logic & \file{chat/peer.py} (\code{ChatNode}), \file{chat/tracker.py}, \file{chat/accounts.py} & peer state, connected peers, message history, registry, users \\
Protocol & \file{chat/protocol.py}, \file{chat/udp_files.py} & framing, encryption, key derivation, UDP reliability and congestion control \\
Transport and HTTP & \file{daemon/} (backend, HTTP adapter, request, response, proxy) & sockets, concurrency modes, HTTP parsing and serialisation \\
\bottomrule
\end{tabularx}
\end{table}

\subsection{Deployment}
On one machine, every process gets its own ports and the browser opens one window per peer. On a LAN,
each peer advertises with \code{--peer-host} the address other machines use to reach it; the tracker
additionally records the address it actually sees (\code{observed_ip}), which makes a wrong
advertised address easy to spot. All peers of a group share the passphrase given with
\code{--shared-key}. The commands are in \cref{app:run}.
